\documentclass[11pt,a4paper]{article}

\usepackage[margin=2cm]{geometry}
\usepackage{tikz}
\usepackage{listings}
\usepackage{xcolor}
\usepackage{hyperref}

\usetikzlibrary{graphs, graphdrawing, shapes.geometric}
\usegdlibrary{trees, layered, force, circular}

\lstset{
  basicstyle=\ttfamily\small,
  breaklines=true,
  frame=single,
  backgroundcolor=\color{black!4},
  columns=fullflexible,
  keepspaces=true,
}

\title{Contoh-Contoh Graph Drawing di TikZ}
\author{Materi dari \url{https://tikz.dev/gd-usage-tikz}}
\date{}

\begin{document}
\maketitle

\section*{Cara Menggunakan Materi Ini}
Setiap contoh berisi: judul, penjelasan singkat, kode sumber, dan hasil gambarnya.
Semua contoh membutuhkan pustaka \texttt{graphs}, \texttt{graphdrawing}, dan algoritma yang dimuat dengan \verb|\usegdlibrary|.
Kompilasi dengan LuaLaTeX.

% ==================================================================
\section{Memilih Layout dan Pustaka (28.1)}

\subsection*{Contoh 1: Layout berlapis dengan panah}
Opsi yang berakhiran \texttt{layout} memanggil algoritma layout. Di sini \texttt{layered layout} menyusun graf dalam lapisan.
\begin{lstlisting}
\begin{tikzpicture}
  \graph [layered layout, edges=rounded corners,
          sibling distance=8mm, level distance=8mm] {
    a -> { b, c -> { d, e } } -> f -> a
  };
\end{tikzpicture}
\end{lstlisting}
\begin{center}
\begin{tikzpicture}
  \graph [layered layout, edges=rounded corners, sibling distance=8mm, level distance=8mm] {
    a -> { b, c -> { d, e } } -> f -> a
  };
\end{tikzpicture}
\end{center}

% ==================================================================
\section{Parameter Graph Drawing (28.2)}

\subsection*{Contoh 2: Segitiga dengan spring layout dan \texttt{vertical}}
\texttt{vertical=1 to 2} membuat garis dari node 1 ke node 2 vertikal. Spring layout mengatur sisanya.
\begin{lstlisting}
\begin{tikzpicture}
  \graph [spring layout, vertical=1 to 2] { 1--2--3--1 };
\end{tikzpicture}
\end{lstlisting}
\begin{center}
\begin{tikzpicture}
  \graph [spring layout, vertical=1 to 2] { 1--2--3--1 };
\end{tikzpicture}
\end{center}

% ==================================================================
\section{Jarak dan Padding (28.3)}

\subsection*{Contoh 3: Jarak node pada necklace (\texttt{node distance})}
Jarak antar node ditentukan oleh \texttt{node distance}. Node 2 dibuat lebih besar sehingga jarak minimal tetap terpenuhi.
\begin{lstlisting}
\begin{tikzpicture}
  \graph [simple necklace layout, node distance=1cm, node sep=0pt,
          nodes={draw, circle, as=.}] {
    1 -- 2[minimum size=2cm] -- 3 -- 4 -- 5 -- 6 -- 7 -- 8[orient=up]
  };
\end{tikzpicture}
\end{lstlisting}
\begin{center}
\begin{tikzpicture}
  \graph [simple necklace layout, node distance=1cm, node sep=0pt,
          nodes={draw, circle, as=.}] {
    1 -- 2[minimum size=2cm] -- 3 -- 4 -- 5 -- 6 -- 7 -- 8[orient=up]
  };
\end{tikzpicture}
\end{center}

\subsection*{Contoh 4: Padding \texttt{node pre sep}}
Padding ditambahkan di depan node. Node 3 dan 5 memakai padding berbeda.
\begin{lstlisting}
\begin{tikzpicture}
  \graph [simple necklace layout, node distance=0cm, node sep=0mm,
          nodes={circle, draw}] {
    1--2--3[node pre sep=5mm]--4--5[node pre sep=1mm]--1
  };
\end{tikzpicture}
\end{lstlisting}
\begin{center}
\begin{tikzpicture}
  \graph [simple necklace layout, node distance=0cm, node sep=0mm,
          nodes={circle, draw}] {
    1--2--3[node pre sep=5mm]--4--5[node pre sep=1mm]--1
  };
\end{tikzpicture}
\end{center}

\subsection*{Contoh 5: \texttt{level distance} dengan \texttt{level pre sep}}
Jarak antar level diatur oleh \texttt{level distance}. Padding per node menambah jarak di sekitarnya.
\begin{lstlisting}
\begin{tikzpicture}
  \graph [layered layout, level distance=0cm, nodes=draw]
    { 1 -- {2,3[level pre sep=1mm],4[level pre sep=5mm]} -- 5[level pre sep=5mm] };
\end{tikzpicture}
\end{lstlisting}
\begin{center}
\begin{tikzpicture}
  \graph [layered layout, level distance=0cm, nodes=draw]
    { 1 -- {2,3[level pre sep=1mm],4[level pre sep=5mm]} -- 5[level pre sep=5mm] };
\end{tikzpicture}
\end{center}

\subsection*{Contoh 6: \texttt{sibling distance} pada pohon}
\texttt{sibling distance=1cm} memberi jarak minimum antar saudara. Dengan \texttt{sibling distance=0cm} dan \texttt{sibling sep=0pt}, saudara akan sangat dekat. Node 3 di-override dengan jarak 1cm.
\begin{lstlisting}
\begin{tikzpicture}
  \graph [tree layout, sibling distance=0cm, sibling sep=0pt,
          nodes={circle, draw}] {
    1 -- {2, 3[sibling distance=1cm], 4, 5}
  };
\end{tikzpicture}
\end{lstlisting}
\begin{center}
\begin{tikzpicture}
  \graph [tree layout, sibling distance=0cm, sibling sep=0pt,
          nodes={circle, draw}] {
    1 -- {2, 3[sibling distance=1cm], 4, 5}
  };
\end{tikzpicture}
\end{center}

\subsection*{Contoh 7: \texttt{sibling pre sep}}
Padding khusus saudara. Node 3 diberi padding di depannya sebesar 1cm.
\begin{lstlisting}
\begin{tikzpicture}
  \graph [tree layout, sibling distance=0cm,
          nodes={circle, draw}, sibling sep=0pt] {
    1 -- {2, 3[sibling pre sep=1cm], 4, 5}
  };
\end{tikzpicture}
\end{lstlisting}
\begin{center}
\begin{tikzpicture}
  \graph [tree layout, sibling distance=0cm,
          nodes={circle, draw}, sibling sep=0pt] {
    1 -- {2, 3[sibling pre sep=1cm], 4, 5}
  };
\end{tikzpicture}
\end{center}

\subsection*{Contoh 8: Komponen dengan \texttt{component sep}}
Dua pohon biner yang tidak terhubung. \texttt{components go right top aligned} meletakkan komponen dari kiri ke kanan dan menyejajarkan bagian atasnya.
\begin{lstlisting}
\begin{tikzpicture}
  \graph [binary tree layout, sibling distance=4mm, level distance=8mm,
          components go right top aligned, component sep=1em,
          nodes=draw] {
    1 -> 2 -> {3 -> 4, 5, 6};
    a -> b -> c -> d -> e;
  };
\end{tikzpicture}
\end{lstlisting}
\begin{center}
\begin{tikzpicture}
  \graph [binary tree layout, sibling distance=4mm, level distance=8mm,
          components go right top aligned, component sep=1em,
          nodes=draw] {
    1 -> 2 -> {3 -> 4, 5, 6};
    a -> b -> c -> d -> e;
  };
\end{tikzpicture}
\end{center}

% ==================================================================
\section{Penjangkaran Graf (28.4)}

\subsection*{Contoh 9: \texttt{desired at} pada spring layout}
Node \texttt{a} diinginkan berada di koordinat (1,2). Algoritma spring akan menggeser seluruh graf agar keinginan ini terpenuhi.
\begin{lstlisting}
\begin{tikzpicture}
  \graph [spring layout] {
    a [desired at={(1,2)}] -- b -- c -- a
  };
\end{tikzpicture}
\end{lstlisting}
\begin{center}
\begin{tikzpicture}
  \graph [spring layout] {
    a [desired at={(1,2)}] -- b -- c -- a
  };
\end{tikzpicture}
\end{center}

\subsection*{Contoh 10: Menggunakan \texttt{x} dan \texttt{y}}
Di dalam graf, kunci \texttt{at} diarahkan ke \texttt{desired at}, sehingga \texttt{x} dan \texttt{y} dapat dipakai langsung.
\begin{lstlisting}
\begin{tikzpicture}
  \graph [layered layout] {
    a [x=1, y=2] -- { b, c } -- { e, f } -- a
  };
\end{tikzpicture}
\end{lstlisting}
\begin{center}
\begin{tikzpicture}
  \graph [layered layout] {
    a [x=1, y=2] -- { b, c } -- { e, f } -- a
  };
\end{tikzpicture}
\end{center}

\subsection*{Contoh 11: Node yang dipaku (\texttt{nail at}) pada spring layout}
Node \texttt{d} dan \texttt{g} dipaku pada posisi tertentu. Node lain bergerak bebas di sekitarnya.
\begin{lstlisting}
\begin{tikzpicture}
  \graph [spring layout] {
    a [x=1] -- { b, c, d [nail at={(0,1)}], e
      -- { f, g [nail at={(2,0)}], h } }
  };
\end{tikzpicture}
\end{lstlisting}
\begin{center}
\begin{tikzpicture}
  \graph [spring layout] {
    a [x=1] -- { b, c, d [nail at={(0,1)}], e
      -- { f, g [nail at={(2,0)}], h } }
  };
\end{tikzpicture}
\end{center}

\subsection*{Contoh 12: \texttt{anchor node} dan \texttt{anchor at}}
Graf dijangkar pada node \texttt{c}, yang diletakkan di koordinat (1,1).
\begin{lstlisting}
\begin{tikzpicture}
  \graph [layered layout, anchor node=c, edges=rounded corners] {
    a -- { b [x=1,y=1], c [x=1,y=1] } -- d -- a
  };
\end{tikzpicture}
\end{lstlisting}
\begin{center}
\begin{tikzpicture}
  \graph [layered layout, anchor node=c, edges=rounded corners] {
    a -- { b [x=1,y=1], c [x=1,y=1] } -- d -- a
  };
\end{tikzpicture}
\end{center}

\subsection*{Contoh 13: \texttt{anchor at} dan \texttt{anchor here}}
Tanpa \texttt{anchor node}, node dengan \texttt{anchor here} menjadi titik jangkar. Node \texttt{c} diletakkan di \texttt{anchor at}, yaitu (1,2).
\begin{lstlisting}
\begin{tikzpicture}
  \graph [layered layout, anchor at={(1,2)}, edges=rounded corners] {
    a -- { b, c [anchor here] } -- d -- a
  };
\end{tikzpicture}
\end{lstlisting}
\begin{center}
\begin{tikzpicture}
  \graph [layered layout, anchor at={(1,2)}, edges=rounded corners] {
    a -- { b, c [anchor here] } -- d -- a
  };
\end{tikzpicture}
\end{center}

% ==================================================================
\section{Orientasi Graf (28.5)}

\subsection*{Contoh 14: Sudut \texttt{orient}}
\texttt{orient=30} membuat garis dari node \texttt{h} ke node \texttt{a} miring 30 derajat.
\begin{lstlisting}
\begin{tikzpicture}
  \graph [spring layout] {
    a -- { b, c, d, e -- {f, g, h} };
    h -- [orient=30] a;
  };
\end{tikzpicture}
\end{lstlisting}
\begin{center}
\begin{tikzpicture}
  \graph [spring layout] {
    a -- { b, c, d, e -- {f, g, h} };
    h -- [orient=30] a;
  };
\end{tikzpicture}
\end{center}

\subsection*{Contoh 15: \texttt{orient'} (membalik graf)}
\texttt{orient'} sama seperti \texttt{orient}, tetapi bagian graf dibalik terhadap garis.
\begin{lstlisting}
\begin{tikzpicture}
  \graph [spring layout] {
    a -- { b, c, d [> orient'=right], e -- {f, g, h} };
    h -- a;
  };
\end{tikzpicture}
\end{lstlisting}
\begin{center}
\begin{tikzpicture}
  \graph [spring layout] {
    a -- { b, c, d [> orient'=right], e -- {f, g, h} };
    h -- a;
  };
\end{tikzpicture}
\end{center}

\subsection*{Contoh 16: \texttt{orient tail} dan \texttt{orient head}}
Garis orientasi dibuat dari node \texttt{f} ke node \texttt{a}, dengan sudut dasar \texttt{orient=|}.
\begin{lstlisting}
\begin{tikzpicture}
  \graph [spring layout] {
    a [orient=|, orient head=f] -- { b, c, d, e -- {f, g, h} };
    { h, g } -- a;
  };
\end{tikzpicture}
\end{lstlisting}
\begin{center}
\begin{tikzpicture}
  \graph [spring layout] {
    a [orient=|, orient head=f] -- { b, c, d, e -- {f, g, h} };
    { h, g } -- a;
  };
\end{tikzpicture}
\end{center}

\subsection*{Contoh 17: \texttt{horizontal} dan \texttt{vertical}}
\texttt{horizontal=a to b} menyejajarkan garis \texttt{a}--\texttt{b} secara horizontal. \texttt{vertical} melakukan hal yang sama secara vertikal.
\begin{lstlisting}
\begin{tikzpicture}
  \graph [spring layout, horizontal=a to b] { a--b--c--a };
  \begin{scope}[xshift=5cm]
    \graph [spring layout, vertical=a to b] { a--b--c--a };
  \end{scope}
\end{tikzpicture}
\end{lstlisting}
\begin{center}
\begin{tikzpicture}
  \graph [spring layout, horizontal=a to b] { a--b--c--a };
  \begin{scope}[xshift=5cm]
    \graph [spring layout, vertical=a to b] { a--b--c--a };
  \end{scope}
\end{tikzpicture}
\end{center}

\subsection*{Contoh 18: \texttt{grow} ke kanan}
\texttt{grow=right} membuat semua anak tumbuh ke kanan. Gunakan \texttt{grow'} jika ingin urutan searah jarum jam.
\begin{lstlisting}
\begin{tikzpicture}
  \graph [layered layout, sibling distance=5mm, grow=right] {
    a -- { b, c, d, e -- {f, g, h} } -- a
  };
\end{tikzpicture}
\end{lstlisting}
\begin{center}
\begin{tikzpicture}
  \graph [layered layout, sibling distance=5mm, grow=right] {
    a -- { b, c, d, e -- {f, g, h} } -- a
  };
\end{tikzpicture}
\end{center}

\subsection*{Contoh 19: \texttt{grow} dengan sudut tertentu}
Graf bipartit lengkap dengan \texttt{grow=-80} agar anak-anak tumbuh pada sudut -80 derajat.
\begin{lstlisting}
\begin{tikzpicture}
  \graph [layered layout, grow=-80] {
    {a,b,c} -- [complete bipartite] {e,d,f}
            -- [complete bipartite] {g,h,i};
  };
\end{tikzpicture}
\end{lstlisting}
\begin{center}
\begin{tikzpicture}
  \graph [layered layout, grow=-80] {
    {a,b,c} -- [complete bipartite] {e,d,f}
            -- [complete bipartite] {g,h,i};
  };
\end{tikzpicture}
\end{center}

% ==================================================================
\section{Menyempurnakan Posisi Node (28.6)}

\subsection*{Contoh 20: \texttt{nudge}}
\texttt{nudge} menggeser node setelah algoritma selesai, sehingga tidak memengaruhi algoritma.
\begin{lstlisting}
\begin{tikzpicture}
  \graph [layered layout, sibling distance=0, nodes=draw, edges=rounded corners] {
    a -- { b, c, d [nudge=(up:2mm)] } -- e -- a
  };
\end{tikzpicture}
\end{lstlisting}
\begin{center}
\begin{tikzpicture}
  \graph [layered layout, sibling distance=0, nodes=draw, edges=rounded corners] {
    a -- { b, c, d [nudge=(up:2mm)] } -- e -- a
  };
\end{tikzpicture}
\end{center}

\subsection*{Contoh 21: \texttt{nudge left} dan \texttt{nudge up}}
Versi singkat dari \texttt{nudge}.
\begin{lstlisting}
\begin{tikzpicture}
  \graph [layered layout, sibling distance=0, nodes=draw, edges=rounded corners] {
    a -- { b [nudge left=2mm], c, d [nudge up=2mm] } -- e -- a
  };
\end{tikzpicture}
\end{lstlisting}
\begin{center}
\begin{tikzpicture}
  \graph [layered layout, sibling distance=0, nodes=draw, edges=rounded corners] {
    a -- { b [nudge left=2mm], c, d [nudge up=2mm] } -- e -- a
  };
\end{tikzpicture}
\end{center}

\subsection*{Contoh 22: \texttt{regardless at}}
Node \texttt{d} dipaksa ke (1,0) setelah algoritma. Algoritma tidak tahu tentang posisi ini.
\begin{lstlisting}
\begin{tikzpicture}
  \graph [layered layout, sibling distance=0, nodes=draw, edges=rounded corners] {
    a -- { b, c, d [regardless at={(1,0)}] } -- e -- a
  };
\end{tikzpicture}
\end{lstlisting}
\begin{center}
\begin{tikzpicture}
  \graph [layered layout, sibling distance=0, nodes=draw, edges=rounded corners] {
    a -- { b, c, d [regardless at={(1,0)}] } -- e -- a
  };
\end{tikzpicture}
\end{center}

\subsection*{Contoh 23: \texttt{nail at}}
Node \texttt{d} dan \texttt{e} ``dipaku'': algoritma diberi tahu posisinya, dan jika gagal, node tetap dipaksa ke sana.
\begin{lstlisting}
\begin{tikzpicture}
  \graph [layered layout, sibling distance=0, nodes=draw, edges=rounded corners] {
    a -- { b, c, d [nail at={(1,0)}] } -- e [nail at={(1.5,-1)}] -- a
  };
\end{tikzpicture}
\end{lstlisting}
\begin{center}
\begin{tikzpicture}
  \graph [layered layout, sibling distance=0, nodes=draw, edges=rounded corners] {
    a -- { b, c, d [nail at={(1,0)}] } -- e [nail at={(1.5,-1)}] -- a
  };
\end{tikzpicture}
\end{center}

% ==================================================================
\section{Packing Komponen Terhubung (28.7)}

\subsection*{Contoh 24: \texttt{componentwise}}
Algoritma necklace biasanya memperlakukan semua node sebagai satu lingkaran. Dengan \texttt{componentwise}, tiap komponen digambar terpisah.
\begin{lstlisting}
\begin{tikzpicture}
  \graph [simple necklace layout, componentwise] {
    a -- b -- c -- d -- a,
    1 -- 2 -- 3 -- 1
  };
\end{tikzpicture}
\end{lstlisting}
\begin{center}
\begin{tikzpicture}
  \graph [simple necklace layout, componentwise] {
    a -- b -- c -- d -- a,
    1 -- 2 -- 3 -- 1
  };
\end{tikzpicture}
\end{center}

\subsection*{Contoh 25: \texttt{component order}}
Urutan komponen menurut kemunculan node pertama dalam input.
\begin{lstlisting}
\begin{tikzpicture}
  \graph [tree layout, nodes={inner sep=1pt, draw, circle},
          component order=by first specified node] {
    a, b, c -- d -- e, f -- g
  };
\end{tikzpicture}
\end{lstlisting}
\begin{center}
\begin{tikzpicture}
  \graph [tree layout, nodes={inner sep=1pt, draw, circle},
          component order=by first specified node] {
    a, b, c -- d -- e, f -- g
  };
\end{tikzpicture}
\end{center}

\subsection*{Contoh 26: \texttt{small components first}}
Komponen diurutkan dari yang terkecil ke terbesar.
\begin{lstlisting}
\begin{tikzpicture}
  \graph [tree layout, nodes={inner sep=1pt, draw, circle},
          small components first] {
    a, b, c -- d -- e, f -- g
  };
\end{tikzpicture}
\end{lstlisting}
\begin{center}
\begin{tikzpicture}
  \graph [tree layout, nodes={inner sep=1pt, draw, circle},
          small components first] {
    a, b, c -- d -- e, f -- g
  };
\end{tikzpicture}
\end{center}

\subsection*{Contoh 27: \texttt{component direction}}
Komponen ditempatkan ke kiri. Kode arah dapat berupa kata (\texttt{left}, \texttt{up}) atau sudut.
\begin{lstlisting}
\begin{tikzpicture}
  \graph [tree layout, nodes={inner sep=1pt, draw, circle},
          component direction=left] {
    a, b, c -- d -- e, f -- g
  };
\end{tikzpicture}
\end{lstlisting}
\begin{center}
\begin{tikzpicture}
  \graph [tree layout, nodes={inner sep=1pt, draw, circle},
          component direction=left] {
    a, b, c -- d -- e, f -- g
  };
\end{tikzpicture}
\end{center}

\subsection*{Contoh 28: \texttt{component align}}
Komponen disejajarkan berdasarkan garis perataan. Nilai \texttt{center} menyejajarkan komponen pada titik tengah proyeksinya.
\begin{lstlisting}
\begin{tikzpicture}
  \graph [binary tree layout, nodes=draw, component align=center] {
    a, b -- c, d -- e -- f
  };
\end{tikzpicture}
\end{lstlisting}
\begin{center}
\begin{tikzpicture}
  \graph [binary tree layout, nodes=draw, component align=center] {
    a, b -- c, d -- e -- f
  };
\end{tikzpicture}
\end{center}

\subsection*{Contoh 29: \texttt{component packing}: rectangular dan skyline}
\texttt{rectangular} membungkus setiap komponen dengan kotak. \texttt{skyline} mendekatkan komponen sejauh mungkin tanpa bersentuhan.
\begin{lstlisting}
\begin{tikzpicture}
  \graph [tree layout, component sep=0pt, component packing=rectangular,
          level distance=1.5cm, nodes=draw] {
    a -- long text, longer text -- b
  };
\end{tikzpicture}
\hspace{1cm}
\begin{tikzpicture}
  \graph [tree layout, component sep=0pt, component packing=skyline,
          level distance=1.5cm, nodes=draw] {
    a -- long text, longer text -- b
  };
\end{tikzpicture}
\end{lstlisting}
\begin{center}
\begin{tikzpicture}
  \graph [tree layout, component sep=0pt, component packing=rectangular,
          level distance=1.5cm, nodes=draw] {
    a -- long text, longer text -- b
  };
\end{tikzpicture}
\hspace{1cm}
\begin{tikzpicture}
  \graph [tree layout, component sep=0pt, component packing=skyline,
          level distance=1.5cm, nodes=draw] {
    a -- long text, longer text -- b
  };
\end{tikzpicture}
\end{center}

% ==================================================================
\section{Menganjurkan Edge (28.8)}

\subsection*{Contoh 30: \texttt{tail anchor} dan \texttt{head anchor}}
Edge dimulai dari sisi timur node awal dan berakhir di sisi barat node akhir.
\begin{lstlisting}
\begin{tikzpicture}
  \graph [layered layout, edges={->}, tail anchor=east,
          head anchor=west, nodes=draw] {
    a -> b -> c
  };
\end{tikzpicture}
\end{lstlisting}
\begin{center}
\begin{tikzpicture}
  \graph [layered layout, edges={->}, tail anchor=east,
          head anchor=west, nodes=draw] {
    a -> b -> c
  };
\end{tikzpicture}
\end{center}

\subsection*{Contoh 31: \texttt{tail cut} dan \texttt{head cut}}
Dengan \texttt{tail cut=false}, edge tidak dipotong di tepi node awal sehingga masuk ke dalam node.
\begin{lstlisting}
\begin{tikzpicture}
  \graph [layered layout, edges={->}, nodes=draw, tail cut=false] {
    a -> b -> c
  };
\end{tikzpicture}
\end{lstlisting}
\begin{center}
\begin{tikzpicture}
  \graph [layered layout, edges={->}, nodes=draw, tail cut=false] {
    a -> b -> c
  };
\end{tikzpicture}
\end{center}

\end{document}
